\documentclass[11pt,a4paper]{article}
\usepackage[T1]{fontenc}
\usepackage{lmodern}
\usepackage[margin=26mm]{geometry}
\usepackage{amsmath,amssymb,microtype}
\setlength{\parindent}{1.2em}
\setlength{\parskip}{0pt}
\newcommand{\sys}{\operatorname{sys}}
\newcommand{\vol}{\operatorname{vol}}
\begin{document}
\begin{center}
{\Large Choosing and changing examples}
\end{center}
\noindent
We used data science to search for convex bodies with large systolic ratios.
One approach selected bodies using geometric quantities that were easier to
compute than capacity. Another changed bodies already sampled, for example by
moving their supporting hyperplanes or rotating them. The experiments below
examine both approaches within the AI-assisted exploration.

For a convex body $K\subset\mathbb R^4$, write
\[
  \sys(K)=\frac{c(K)^2}{2\vol(K)},
\]
where $c(K)$ is its Ekeland--Hofer--Zehnder capacity and $\vol$ denotes
four-dimensional volume. The search sought values above $1$. Multiplying all
coordinates by $t>0$ multiplies capacity by $t^2$ and volume by $t^4$, so the
ratio does not change. Simply making a body larger cannot help.

Many of the examples were products of two planar polygons. In coordinates
$(q_1,q_2,p_1,p_2)$, we formed
\[
 K=P\times Q,\qquad
 P=\{q\in\mathbb R^2:\langle n_i,q\rangle\le h_i,\ i=1,\ldots,k\},
\]
with an analogous description of $Q$ using $m$ inequalities. The unit normals
$n_i$ came from independent uniform angles, and the heights $h_i$ were sampled
independently from $[0.8,1.2)$. We retained bounded polygons for which every
proposed side was present. Thus a $k\times m$ sample means a product with $k$
sides in the first factor and $m$ in the second; the rejection rule is part of
its distribution. The computations below report numerical ratios, obtained by
substituting computed capacity and volume into the displayed formula. The new
experiments bound the error in capacity, but their volume calculation still uses
floating-point arithmetic.

A useful descriptor emerged from the two-dimensional faces of these bodies.
Such a face is called a ridge in dimension four. If a ridge $R$ has cyclically
ordered vertices $v_0,\ldots,v_{r-1}$, its absolute symplectic area is
\[
 A_\omega(R)=\frac12\left|\sum_{i=0}^{r-1}\omega_0(v_i,v_{i+1})\right|,
 \qquad v_r=v_0,
 \qquad \omega_0=\sum_{j=1}^2dq_j\wedge dp_j.
\]
This is the polygonal area formula with the symplectic pairing in place of the
planar determinant. It measures the area seen by $\omega_0$, which can vanish
even when the face has positive Euclidean area. Summing over the ridges gives
the descriptor
\[
 S(K)=\frac{\sum_R A_\omega(R)}{\sqrt{\vol(K)}}.
\]
Both numerator and denominator scale by $t^2$, so $S$ is unchanged by a common
dilation. It is also unchanged by translation and linear symplectic maps.

In the original samples, smaller values of $S$ tended to accompany larger
ratios. The association persisted when products with the same numbers of sides
were considered separately. This suggested a simple use: generate many bodies,
compute their ridge areas, and select those with low $S$ before calculating
their capacities. Tests on fresh candidates showed that low-ridge selection
could raise the ratios relative to controls. That success supplied a way to
choose examples. It did not yet tell us which change to make to a given body.

One concrete change is to replace all support heights of a polygon by $1$,
keeping its normals fixed. Every resulting side lies on a tangent line to the
unit circle. We tested this operation on sixteen independently generated pairs
of polygons: eight of type $4\times4$ and eight of type $4\times6$. Each pair
supplied four bodies: the original product, the product with only the first
factor changed, the product with only the second changed, and the product with
both changed. We accepted a pair only if all four constructions were valid,
before evaluating any of their ratios. The comparison therefore concerns this
jointly admissible sample.

Each planar factor was then scaled to area one. This normalization does not
alter the mathematical ratio: separate factor dilations have the form
$\operatorname{diag}(aI_2,bI_2)$, which is a common dilation composed with the
symplectic map
$\operatorname{diag}(\sqrt{a/b}\,I_2,\sqrt{b/a}\,I_2)$.
The shape change comes from replacing unequal heights, not from adjusting the
areas. Changing both factors increased the mean numerical ratio by $0.0153$.
Eleven pairs improved and five became worse. The mean increase was about
$0.004$ for $4\times4$ and $0.027$ for $4\times6$. These paired comparisons
show why an average gain should not become a rule that the operation always
helps: even with the same normals held fixed, its effect varied across bodies.

Rotation asks a different question. An orthogonal map preserves Euclidean shape,
but it need not preserve $\omega_0$. Rotating a Lagrangian product can therefore
change its capacity; the rotated body need no longer be a Lagrangian product
for the fixed symplectic form. This gives a way to explore symplectic alignment
without changing the body's Euclidean geometry.

To test whether that freedom was worth using in a search, we compared two
allocations of sixteen evaluations. One used a single fresh product, one
symplectic rotation as a control, and fourteen prescribed orthogonal rotations.
The other evaluated sixteen fresh products, sharing the first body with the
rotation experiment. We made this comparison four times, twice for $3\times3$
and twice for $4\times4$. The rotations and all source bodies were fixed before
their ratios were evaluated; the symplectic controls preserved the ratios to
within $4\cdot10^{-16}$.

Rotation improved three of the four starting bodies. Nevertheless, fresh
sampling produced the larger maximum in every comparison. In one $4\times4$
case, rotation raised the ratio from $0.328$ to $0.529$, while the sixteen fresh
products reached $0.774$. The rotation calculations also took longer. For these four starting bodies, spending sixteen evaluations on fresh products
was more effective than spending them on the prescribed rotations.
\end{document}
